% ECONOMICS_PROBLEM: {"id": "C65-26", "title": "Quadratic BSDE uniqueness with only one exponential moment", "jel_code": "C65", "source_id": "OP-0565"}
% !TeX program = xelatex
\documentclass[11pt,a4paper]{article}
\usepackage[margin=23mm]{geometry}
\usepackage{fontspec}
\setmainfont{Latin Modern Roman}
\usepackage{amsmath,amssymb,mathtools}
\usepackage{xcolor,enumitem,booktabs,longtable,array,xurl,hyperref}
\definecolor{muted}{HTML}{53636C}
\hypersetup{colorlinks=true,urlcolor=blue,linkcolor=blue}
\setlength{\parindent}{0pt}
\setlength{\parskip}{5pt}
\newcommand{\R}{\mathbb R}
\newcommand{\E}{\mathbb E}
\newcommand{\N}{\mathbb N}
\newcommand{\ind}{\mathbf 1}
\newcommand{\Var}{\operatorname{Var}}
\newcommand{\Cov}{\operatorname{Cov}}
\newcommand{\supp}{\operatorname{supp}}
\DeclareMathOperator*{\argmin}{arg\,min}
\DeclareMathOperator*{\argmax}{arg\,max}
\newcommand{\entrymeta}[3]{{\sffamily\small\color{muted}\textbf{\texttt{#1}}\quad|\quad #2\par #3\par}}
\newcommand{\Problemref}[1]{\texttt{#1}}
\newcommand{\JELref}[2]{\texttt{#2} (JEL \texttt{#1})}
\begin{document}
\section*{Quadratic BSDE uniqueness with only one exponential moment}
\textbf{Problem ID:} \texttt{C65-26}\quad \textbf{JEL:} \texttt{C65}\par
% BEGIN ECONOMICS STATEMENT
\label{problem:OP-0565}
\label{beyond:SC-04}
\entrymeta{SC-04}{Stochastic control and dynamic optimization}{Published question}
\subsection*{Definitions and assumptions}
Use the scalar Brownian BSDE conventions above. Let
$g:\R^d\to\R$ be a finite convex function satisfying
\[
 0\leq g(z)\leq\tfrac12|z|^2\qquad(z\in\R^d).
\]
Convexity means
$g(\theta z+(1-\theta)z')\leq\theta g(z)+(1-\theta)g(z')$
for $\theta\in[0,1]$; no positive lower bound on the curvature is
assumed. Let $\xi$ satisfy $\E e^{|\xi|}<\infty$.
The relevant solution class consists of scalar BSDE solutions for which
$(e^{|Y_t|})_{0\leq t\leq T}$ is of class (D).
\subsection*{Formal question}
For every generator and terminal value satisfying these hypotheses,
does the equation
\[
 Y_t=\xi+\int_t^T g(Z_s)\,ds-\int_t^T Z_s\cdot dB_s
\]
have at most one solution in this exponential class-(D) class?
Equivalently, must any two solutions in this class have indistinguishable
$Y$ components and equal $Z$ components $dt\otimes\Pr$ almost everywhere?
\subsection*{Scope and status}
The coefficient $1/2$ in the quadratic upper bound and the exponent
$1$ in both integrability requirements fix the normalization of the
question. A theorem requiring $\E e^{q|\xi|}<\infty$ and class (D)
for $e^{q|Y|}$ with $q>1$ would address a stronger-data regime.
A theorem requiring strong convexity would address a smaller generator
class. Neither modification answers this endpoint formulation. The
question is explicitly recorded in the cited source version; no
unverified assertion of current universal openness is made here.
\subsection*{References for the source of the problem}
Fan, Hu and Tang, \emph{1D nonlinear backward stochastic differential
equations: a unified theory and applications}, arXiv:2309.06233v2,
\S5, Problem 5.4 and equation (5.3). That paragraph identifies the
strong-convexity and stronger-exponential-moment results against which
the endpoint question is posed.
\url{https://arxiv.org/html/2309.06233v2#S5}.

\subsection*{Common conventions: Source mathematical conventions}

\textbf{Well-defined objects.}
All probability spaces and measurable variables are specified or inherited
from a local statistical model. Expectations used as finite targets require
the stated integrability. A decision rule or estimator is measurable with
respect to its available information; randomized rules may use an auxiliary
independent random variable. Norms without a subscript are Euclidean in
finite-dimensional spaces. An optimizer is requested only when attainment
is assured; otherwise the question concerns the optimal value and
$\varepsilon$-optimal admissible rules for every $\varepsilon>0$.

\textbf{Identification.}
A structural model $m\in\mathcal M$ implies an observable law
$P_m^{\rm obs}$ and target $\theta(m)$. The target is point identified on
$\mathcal M$ if
\[
 P_{m_1}^{\rm obs}=P_{m_2}^{\rm obs}
 \quad\Longrightarrow\quad\theta(m_1)=\theta(m_2)
 \qquad(m_1,m_2\in\mathcal M).
\]
When this fails, the sharp identified set is
\[
 \Theta(P;\mathcal M)=\{\theta(m):m\in\mathcal M,
                                  P_m^{\rm obs}=P\}.
\]
Specifying an intervention or estimand does not establish identification.
Positivity, exclusion, causal sufficiency, measurement restrictions, and
transport assumptions are substantive local inputs, never automatic
consequences of notation.

\textbf{Minimax risk and nontrivial inference.}
For a statistical experiment with data law $P^{(n)}$, target $\theta(P)$,
and nonnegative loss $\ell$, define
\[
 \mathcal R_n(\mathcal P,\ell)
   =\inf_{\widehat\theta_n}\sup_{P\in\mathcal P}
      \E_{P^{(n)}}\ell(\widehat\theta_n,\theta(P)).
\]
The infimum is over measurable estimators. A sharp rate is a sequence
$r_n$ for which $0<c\leq C<\infty$, independent of $n$, satisfies
$c r_n\leq\mathcal R_n\leq C r_n$ for all sufficiently large $n$.
Squared-error parametric risk is of order $n^{-1}$; the corresponding
root-mean-squared error is of order $n^{-1/2}$. These different conventions
are distinguished locally. A uniform asymptotic confidence region $C_n$
must satisfy
\[
 \liminf_{n\to\infty}\inf_{P\in\mathcal P}
     \Pr_{P^{(n)}}\{\theta(P)\in C_n\}\geq 1-\alpha,
 \qquad 0<\alpha<1,
\]
together with an informative diameter or expected-width requirement.
Coverage alone permits the vacuous whole parameter space. Testing questions
similarly impose both size control and power against a declared separated
alternative class.

\textbf{Binary causal baseline.}
When an entry explicitly invokes this baseline, one observes independent
copies of $O=(X,W,Y)$, with $W\in\{0,1\}$ and potential outcomes $Y(0),Y(1)$.
Assume consistency $Y=Y(W)$, no interference, conditional exchangeability
$(Y(0),Y(1))\perp W\mid X$, and overlap
$\epsilon\leq e(X)\leq1-\epsilon$ almost surely for a fixed
$0<\epsilon<1/2$, where
\[
 e(x)=\Pr(W=1\mid X=x),\qquad
 m_w(x)=\E[Y\mid W=w,X=x],\qquad
 \tau(x)=m_1(x)-m_0(x).
\]
These assumptions identify conditional mean effects almost everywhere
under the covariate law. Pointwise or uniform effects require appropriate
regularity and a specified support. Continuous treatment, longitudinal
data, adaptive experiments, and network interference require their own
assumptions in place of this baseline. Bounded outcomes, densities, and
smoothness are additional assumptions only where stated.

\textbf{Function classes.}
On $[0,1]^d$, $\mathcal H_d^s(L)$ is the isotropic Hölder ball of radius
$L>0$: put $k=\lceil s\rceil-1$ and $\gamma=s-k\in(0,1]$, bound derivatives
through order $k$, and require the order-$k$ derivatives to be
$\gamma$-Hölder with constant at most $L$. Thus integer orders use a
Lipschitz final derivative convention. The class
$\operatorname{BV}_d(L)$ consists of functions $f\in L^1((0,1)^d)$ whose
distributional derivative is a finite vector measure and for which
$\|f\|_{L^1}+|Df|((0,1)^d)\leq L$.
An $L^\infty$ bound or a particular pointwise version is imposed separately
when needed. Sparse and neural-network classes require a declared
dictionary or architecture and coefficient bounds.

An anisotropic H\"older class with declared block smoothness indices means
coordinatewise H\"older regularity in each block, uniformly in the other
blocks, with all corresponding derivative and H\"older bounds at most
the stated radius. Derivative orders and the integer-order convention in
each block are those just given. Mixed-derivative bounds are additional
restrictions only when explicitly stated.

\textbf{Decision and computational questions.}
Policy regret compares admissible feasible policies under the same law and
constraints. In computational entries, finite data and thresholds have
rational encodings; running time is measured in their total bit length.
Real-valued equilibrium search uses the explicitly declared algebraic or
FIXP encoding. An approximation ratio for maximization is written
$\operatorname{OPT}/\operatorname{ALG}$ and for minimization as
$\operatorname{ALG}/\operatorname{OPT}$, with zero optima and infeasible
instances handled locally. Historical approximation bounds are not
automatically current bounds.

\subsection*{Common conventions: BSDE conventions for SC-01--SC-06}

Fix $T>0$ and an integer $d\geq1$. The probability space
$(\Omega,\mathcal F,\Pr)$ is complete, $B$ is a standard $d$-dimensional
Brownian motion, and $(\mathcal F_t)_{0\leq t\leq T}$ is its usual augmented
natural filtration, with $\mathcal F=\mathcal F_T$. A scalar BSDE solution
is a pair $(Y,Z)$, where $Y$ is continuous and adapted, $Z$ is predictable
and $\R^d$-valued, and, almost surely,
\[
 \int_0^T\bigl(|g(s,Y_s,Z_s)|+|Z_s|^2\bigr)\,ds<\infty,
 \qquad
 Y_t=\xi+\int_t^T g(s,Y_s,Z_s)\,ds-\int_t^T Z_s\cdot dB_s.
\]
The identity holds simultaneously for every $t\in[0,T]$. Uniqueness means
indistinguishability of $Y$ and equality of $Z$ for $dt\otimes\Pr$ almost
every $(t,\omega)$. A process $U$ is of class (D) when
$\{U_\tau:\tau\text{ is an }(\mathcal F_t)\text{-stopping time},\ 0\leq\tau\leq T\}$
is uniformly integrable. For $p>0$, $L^p$ denotes
$\{\xi:\E|\xi|^p<\infty\}$; for $\mu>0$,
$\exp(\mu L)$ denotes $\{\xi:\E e^{\mu|\xi|}<\infty\}$.
These conventions remove ambiguity about filtration, stochastic integration,
and integrability; they do not add an $L^2$ expectation requirement to the
scalar solution definition. The conventions follow Fan, Hu and Tang,
\emph{1D nonlinear backward stochastic differential equations: a unified
theory and applications}, arXiv:2309.06233v2 (11 February 2026),
\S\S1 and 1.3, \url{https://arxiv.org/html/2309.06233v2}.
% END ECONOMICS STATEMENT
\end{document}
